% --- Archon physics formalization source begin ---
% archon:physics
% archon:covers PhyXMiniProblems/problem_phyx_mini_0747.lean
% archon:source-report reports/phyx_mini/problem_phyx_mini_0747.source.json

\chapter{Physics problem phyx\_mini\_0747}
\label{ch:PhyXMiniProblems_problem_phyx_mini_0747}

\paragraph{Problem source.}
\#\# Physical scenario

In figure, a baseball is hit at a height \$h = 1.00\textbackslash{} m\$ and then caught at the same height. It travels alongside a wall, moving up past the top of the wall \$1.00\textbackslash{} s\$ after it is hit and then down past the top of the wall \$4.00\textbackslash{} s\$ later, at distance \$D = 50.0\textbackslash{} m\$ farther along the wall.

\#\# Question

How high is the wall?

\#\# Answer choices

- A: A: 21.5 m

- B: B: 23.5 m

- C: C: 25.5 m

- D: D: 27.5 m

\#\# Figure caption (auxiliary; use the image as primary evidence)

The image shows a diagram of a brick wall, topped with a dashed blue parabolic arc representing a trajectory. 

- The wall is made of uniformly-sized rectangular bricks, stacked to form a solid structure.
- The parabolic arc stretches over the top of the wall, illustrating a movement or trajectory.
- Two vertical arrows near each end of the arc indicate a height labeled \textbackslash{}( h \textbackslash{}).
- A horizontal double-headed arrow beneath the arc represents a distance labeled \textbackslash{}( D \textbackslash{}). 

The diagram seems to illustrate the physics of a projectile motion, such as an object launched over a barrier like a wall.

\#\# Dataset metadata

Kinematics; Multi-Formula Reasoning, Physical Model Grounding Reasoning

\paragraph{Recorded answer/context.}
C: C: 25.5 m

\paragraph{Figure/image path.}
/root/proposal\_for\_physic/science-mango/phyx\_mini\_run/phyx\_data/test\_image/747.png

\paragraph{Formalization target.}
create a compiling Lean file with sorry bodies at `PhyXMiniProblems/problem\_phyx\_mini\_0747.lean`. The Lean declarations must preserve the physical quantities, dimensions or dimensional roles, figure labels, governing-law hypotheses, and final relation expressed by this problem.
Use Mathlib/Physlib names found through LeanExplore where available. If a physics API is missing, introduce faithful local abstractions rather than scalar placeholder aliases.

\begin{theorem}[Physics formalization target]
\label{thm:physics:phyx_mini_0747:target}
\lean{PhyXMiniProblems.ProblemPhyXMini0747.problem_phyx_mini_0747}
\uses{def:physics:phyx-mini-0747:phyxminiproblems-problemphyxmini0747-lengthinmeters, def:physics:phyx-mini-0747:phyxminiproblems-problemphyxmini0747-baseballwallsetup, def:physics:phyx-mini-0747:phyxminiproblems-problemphyxmini0747-matchesproblemstatement, def:physics:phyx-mini-0747:phyxminiproblems-problemphyxmini0747-matchesprimaryfigure, def:physics:phyx-mini-0747:phyxminiproblems-problemphyxmini0747-usesstandardnearearthgravity, def:physics:phyx-mini-0747:phyxminiproblems-problemphyxmini0747-hasphysicalprojectileparameters, def:physics:phyx-mini-0747:phyxminiproblems-problemphyxmini0747-satisfiesidealprojectilelaws, def:physics:phyx-mini-0747:phyxminiproblems-problemphyxmini0747-matchesdisplayedwallheight}
The assigned autoformalize agent should translate this physics problem into Lean declarations in the covered file, with theorem and lemma proof bodies written as `by sorry`.
\end{theorem}
\begin{proof}
This is an autoformalization task, not a proof task. Produce statements that can later be proved by the physics prover without weakening the physical model.
\end{proof}

% --- Archon declaration topology begin ---
\section{Declaration topology}
\begin{definition}[velocity Dimension]
  \label{def:physics:phyx-mini-0747:phyxminiproblems-problemphyxmini0747-velocitydimension}
  \lean{PhyXMiniProblems.ProblemPhyXMini0747.velocityDimension}
  The physical dimension of velocity, L T\textasciicircum{}-1
\end{definition}
\begin{proof}
  This is the declared model definition of velocity Dimension.
\end{proof}
\begin{definition}[acceleration Dimension]
  \label{def:physics:phyx-mini-0747:phyxminiproblems-problemphyxmini0747-accelerationdimension}
  \lean{PhyXMiniProblems.ProblemPhyXMini0747.accelerationDimension}
  The physical dimension of acceleration, L T\textasciicircum{}-2
\end{definition}
\begin{proof}
  This is the declared model definition of acceleration Dimension.
\end{proof}
\begin{definition}[Length Quantity]
  \label{def:physics:phyx-mini-0747:phyxminiproblems-problemphyxmini0747-lengthquantity}
  \lean{PhyXMiniProblems.ProblemPhyXMini0747.LengthQuantity}
  A nonnegative physical length, independent of the unit used to read it
\end{definition}
\begin{proof}
  This is the declared model definition of Length Quantity.
\end{proof}
\begin{definition}[Signed Length Quantity]
  \label{def:physics:phyx-mini-0747:phyxminiproblems-problemphyxmini0747-signedlengthquantity}
  \lean{PhyXMiniProblems.ProblemPhyXMini0747.SignedLengthQuantity}
  A signed Cartesian position component with the dimension of length
\end{definition}
\begin{proof}
  This is the declared model definition of Signed Length Quantity.
\end{proof}
\begin{definition}[Time Quantity]
  \label{def:physics:phyx-mini-0747:phyxminiproblems-problemphyxmini0747-timequantity}
  \lean{PhyXMiniProblems.ProblemPhyXMini0747.TimeQuantity}
  A nonnegative physical elapsed time
\end{definition}
\begin{proof}
  This is the declared model definition of Time Quantity.
\end{proof}
\begin{definition}[Signed Velocity Quantity]
  \label{def:physics:phyx-mini-0747:phyxminiproblems-problemphyxmini0747-signedvelocityquantity}
  \lean{PhyXMiniProblems.ProblemPhyXMini0747.SignedVelocityQuantity}
  \uses{def:physics:phyx-mini-0747:phyxminiproblems-problemphyxmini0747-velocitydimension}
  A signed Cartesian velocity component
\end{definition}
\begin{proof}
  This is the declared model definition of Signed Velocity Quantity.
\end{proof}
\begin{definition}[Acceleration Magnitude Quantity]
  \label{def:physics:phyx-mini-0747:phyxminiproblems-problemphyxmini0747-accelerationmagnitudequantity}
  \lean{PhyXMiniProblems.ProblemPhyXMini0747.AccelerationMagnitudeQuantity}
  \uses{def:physics:phyx-mini-0747:phyxminiproblems-problemphyxmini0747-accelerationdimension}
  A nonnegative magnitude of gravitational acceleration
\end{definition}
\begin{proof}
  This is the declared model definition of Acceleration Magnitude Quantity.
\end{proof}
\begin{definition}[length In Meters]
  \label{def:physics:phyx-mini-0747:phyxminiproblems-problemphyxmini0747-lengthinmeters}
  \lean{PhyXMiniProblems.ProblemPhyXMini0747.lengthInMeters}
  \uses{def:physics:phyx-mini-0747:phyxminiproblems-problemphyxmini0747-lengthquantity}
  Read a nonnegative physical length in SI metres
\end{definition}
\begin{proof}
  This is the declared model definition of length In Meters.
\end{proof}
\begin{definition}[signed Length In Meters]
  \label{def:physics:phyx-mini-0747:phyxminiproblems-problemphyxmini0747-signedlengthinmeters}
  \lean{PhyXMiniProblems.ProblemPhyXMini0747.signedLengthInMeters}
  \uses{def:physics:phyx-mini-0747:phyxminiproblems-problemphyxmini0747-signedlengthquantity}
  Read a signed Cartesian position component in SI metres
\end{definition}
\begin{proof}
  This is the declared model definition of signed Length In Meters.
\end{proof}
\begin{definition}[time In Seconds]
  \label{def:physics:phyx-mini-0747:phyxminiproblems-problemphyxmini0747-timeinseconds}
  \lean{PhyXMiniProblems.ProblemPhyXMini0747.timeInSeconds}
  \uses{def:physics:phyx-mini-0747:phyxminiproblems-problemphyxmini0747-timequantity}
  Read a physical elapsed time in SI seconds
\end{definition}
\begin{proof}
  This is the declared model definition of time In Seconds.
\end{proof}
\begin{definition}[signed Velocity In Meters Per Second]
  \label{def:physics:phyx-mini-0747:phyxminiproblems-problemphyxmini0747-signedvelocityinmeterspersecond}
  \lean{PhyXMiniProblems.ProblemPhyXMini0747.signedVelocityInMetersPerSecond}
  \uses{def:physics:phyx-mini-0747:phyxminiproblems-problemphyxmini0747-signedvelocityquantity}
  Read a signed Cartesian velocity component in SI metres per second
\end{definition}
\begin{proof}
  This is the declared model definition of signed Velocity In Meters Per Second.
\end{proof}
\begin{definition}[acceleration In Meters Per Second Squared]
  \label{def:physics:phyx-mini-0747:phyxminiproblems-problemphyxmini0747-accelerationinmeterspersecondsquared}
  \lean{PhyXMiniProblems.ProblemPhyXMini0747.accelerationInMetersPerSecondSquared}
  \uses{def:physics:phyx-mini-0747:phyxminiproblems-problemphyxmini0747-accelerationmagnitudequantity}
  Read an acceleration magnitude in SI metres per second squared
\end{definition}
\begin{proof}
  This is the declared model definition of acceleration In Meters Per Second Squared.
\end{proof}
\begin{definition}[Projectile Kind]
  \label{def:physics:phyx-mini-0747:phyxminiproblems-problemphyxmini0747-projectilekind}
  \lean{PhyXMiniProblems.ProblemPhyXMini0747.ProjectileKind}
  The type of moving object named in the problem
\end{definition}
\begin{proof}
  This is the declared model definition of Projectile Kind.
\end{proof}
\begin{definition}[Projectile Model]
  \label{def:physics:phyx-mini-0747:phyxminiproblems-problemphyxmini0747-projectilemodel}
  \lean{PhyXMiniProblems.ProblemPhyXMini0747.ProjectileModel}
  The idealized mechanics model used after the hit
\end{definition}
\begin{proof}
  This is the declared model definition of Projectile Model.
\end{proof}
\begin{definition}[Height Endpoint]
  \label{def:physics:phyx-mini-0747:phyxminiproblems-problemphyxmini0747-heightendpoint}
  \lean{PhyXMiniProblems.ProblemPhyXMini0747.HeightEndpoint}
  The hit and catch endpoints whose common height is labelled h
\end{definition}
\begin{proof}
  This is the declared model definition of Height Endpoint.
\end{proof}
\begin{definition}[Wall Top Crossing]
  \label{def:physics:phyx-mini-0747:phyxminiproblems-problemphyxmini0747-walltopcrossing}
  \lean{PhyXMiniProblems.ProblemPhyXMini0747.WallTopCrossing}
  The two intersections of the dashed trajectory with the wall top
\end{definition}
\begin{proof}
  This is the declared model definition of Wall Top Crossing.
\end{proof}
\begin{definition}[Figure Label]
  \label{def:physics:phyx-mini-0747:phyxminiproblems-problemphyxmini0747-figurelabel}
  \lean{PhyXMiniProblems.ProblemPhyXMini0747.FigureLabel}
  Literal mathematical labels visible in image 747.png
\end{definition}
\begin{proof}
  This is the declared model definition of Figure Label.
\end{proof}
\begin{definition}[Figure Object]
  \label{def:physics:phyx-mini-0747:phyxminiproblems-problemphyxmini0747-figureobject}
  \lean{PhyXMiniProblems.ProblemPhyXMini0747.FigureObject}
  Physical or graphical objects visibly represented in the primary image
\end{definition}
\begin{proof}
  This is the declared model definition of Figure Object.
\end{proof}
\begin{definition}[Figure Anchor]
  \label{def:physics:phyx-mini-0747:phyxminiproblems-problemphyxmini0747-figureanchor}
  \lean{PhyXMiniProblems.ProblemPhyXMini0747.FigureAnchor}
  Schematic anchors used as endpoints of the three dimension arrows
\end{definition}
\begin{proof}
  This is the declared model definition of Figure Anchor.
\end{proof}
\begin{definition}[Planar Position]
  \label{def:physics:phyx-mini-0747:phyxminiproblems-problemphyxmini0747-planarposition}
  \lean{PhyXMiniProblems.ProblemPhyXMini0747.PlanarPosition}
  \uses{def:physics:phyx-mini-0747:phyxminiproblems-problemphyxmini0747-signedlengthquantity}
  A planar physical position with signed horizontal and vertical components
\end{definition}
\begin{proof}
  This is the declared model definition of Planar Position.
\end{proof}
\begin{definition}[Planar Velocity]
  \label{def:physics:phyx-mini-0747:phyxminiproblems-problemphyxmini0747-planarvelocity}
  \lean{PhyXMiniProblems.ProblemPhyXMini0747.PlanarVelocity}
  \uses{def:physics:phyx-mini-0747:phyxminiproblems-problemphyxmini0747-signedvelocityquantity}
  A planar physical velocity with signed horizontal and vertical components
\end{definition}
\begin{proof}
  This is the declared model definition of Planar Velocity.
\end{proof}
\begin{definition}[Baseball Wall Figure]
  \label{def:physics:phyx-mini-0747:phyxminiproblems-problemphyxmini0747-baseballwallfigure}
  \lean{PhyXMiniProblems.ProblemPhyXMini0747.BaseballWallFigure}
  \uses{def:physics:phyx-mini-0747:phyxminiproblems-problemphyxmini0747-lengthquantity, def:physics:phyx-mini-0747:phyxminiproblems-problemphyxmini0747-heightendpoint, def:physics:phyx-mini-0747:phyxminiproblems-problemphyxmini0747-walltopcrossing, def:physics:phyx-mini-0747:phyxminiproblems-problemphyxmini0747-figurelabel, def:physics:phyx-mini-0747:phyxminiproblems-problemphyxmini0747-figureobject, def:physics:phyx-mini-0747:phyxminiproblems-problemphyxmini0747-figureanchor}
  The labels, dimension-arrow endpoints, and qualitative geometry read directly from the supplied bitmap. The marked quantities are independent physical lengths; the figure does not assign a numerical value to the wall height
\end{definition}
\begin{proof}
  This is the declared model definition of Baseball Wall Figure.
\end{proof}
\begin{definition}[Baseball Wall Setup]
  \label{def:physics:phyx-mini-0747:phyxminiproblems-problemphyxmini0747-baseballwallsetup}
  \lean{PhyXMiniProblems.ProblemPhyXMini0747.BaseballWallSetup}
  \uses{def:physics:phyx-mini-0747:phyxminiproblems-problemphyxmini0747-lengthquantity, def:physics:phyx-mini-0747:phyxminiproblems-problemphyxmini0747-timequantity, def:physics:phyx-mini-0747:phyxminiproblems-problemphyxmini0747-signedvelocityquantity, def:physics:phyx-mini-0747:phyxminiproblems-problemphyxmini0747-accelerationmagnitudequantity, def:physics:phyx-mini-0747:phyxminiproblems-problemphyxmini0747-projectilekind, def:physics:phyx-mini-0747:phyxminiproblems-problemphyxmini0747-projectilemodel, def:physics:phyx-mini-0747:phyxminiproblems-problemphyxmini0747-planarposition, def:physics:phyx-mini-0747:phyxminiproblems-problemphyxmini0747-planarvelocity, def:physics:phyx-mini-0747:phyxminiproblems-problemphyxmini0747-baseballwallfigure}
  All physical quantities and observables in the experiment. The wall height is an independent field, not a definition of the requested answer. The position and velocity functions are constrained only by the governing laws below
\end{definition}
\begin{proof}
  This is the declared model definition of Baseball Wall Setup.
\end{proof}
\begin{definition}[Matches Problem Statement]
  \label{def:physics:phyx-mini-0747:phyxminiproblems-problemphyxmini0747-matchesproblemstatement}
  \lean{PhyXMiniProblems.ProblemPhyXMini0747.MatchesProblemStatement}
  \uses{def:physics:phyx-mini-0747:phyxminiproblems-problemphyxmini0747-lengthinmeters, def:physics:phyx-mini-0747:phyxminiproblems-problemphyxmini0747-signedlengthinmeters, def:physics:phyx-mini-0747:phyxminiproblems-problemphyxmini0747-timeinseconds, def:physics:phyx-mini-0747:phyxminiproblems-problemphyxmini0747-baseballwallsetup}
  Numerical and incidence information supplied by the prose. In particular, the two trajectory values at the crossing times are equated to the independent wall-height field, but that field is not assigned the requested numerical answer here
\end{definition}
\begin{proof}
  This is the declared model definition of Matches Problem Statement.
\end{proof}
\begin{definition}[Matches Primary Figure]
  \label{def:physics:phyx-mini-0747:phyxminiproblems-problemphyxmini0747-matchesprimaryfigure}
  \lean{PhyXMiniProblems.ProblemPhyXMini0747.MatchesPrimaryFigure}
  \uses{def:physics:phyx-mini-0747:phyxminiproblems-problemphyxmini0747-baseballwallsetup}
  Literal and qualitative evidence from image 747.png. These fields connect the repeated h label to the common hit/catch height and the D label to the distance between wall-top crossings, without asserting the wall's height
\end{definition}
\begin{proof}
  This is the declared model definition of Matches Primary Figure.
\end{proof}
\begin{definition}[Uses Standard Near Earth Gravity]
  \label{def:physics:phyx-mini-0747:phyxminiproblems-problemphyxmini0747-usesstandardnearearthgravity}
  \lean{PhyXMiniProblems.ProblemPhyXMini0747.UsesStandardNearEarthGravity}
  \uses{def:physics:phyx-mini-0747:phyxminiproblems-problemphyxmini0747-accelerationinmeterspersecondsquared, def:physics:phyx-mini-0747:phyxminiproblems-problemphyxmini0747-baseballwallsetup}
  The standard near-Earth gravitational acceleration implicit in the recorded numerical answer. This calibration contains no wall-height information
\end{definition}
\begin{proof}
  This is the declared model definition of Uses Standard Near Earth Gravity.
\end{proof}
\begin{definition}[Has Physical Projectile Parameters]
  \label{def:physics:phyx-mini-0747:phyxminiproblems-problemphyxmini0747-hasphysicalprojectileparameters}
  \lean{PhyXMiniProblems.ProblemPhyXMini0747.HasPhysicalProjectileParameters}
  \uses{def:physics:phyx-mini-0747:phyxminiproblems-problemphyxmini0747-lengthinmeters, def:physics:phyx-mini-0747:phyxminiproblems-problemphyxmini0747-timeinseconds, def:physics:phyx-mini-0747:phyxminiproblems-problemphyxmini0747-signedvelocityinmeterspersecond, def:physics:phyx-mini-0747:phyxminiproblems-problemphyxmini0747-accelerationinmeterspersecondsquared, def:physics:phyx-mini-0747:phyxminiproblems-problemphyxmini0747-baseballwallsetup}
  Positivity, chronological order, and the ascending/descending branches named in the prose. These select the physical trajectory without fixing the wall height to any displayed answer
\end{definition}
\begin{proof}
  This is the declared model definition of Has Physical Projectile Parameters.
\end{proof}
\begin{definition}[Satisfies Ideal Projectile Laws]
  \label{def:physics:phyx-mini-0747:phyxminiproblems-problemphyxmini0747-satisfiesidealprojectilelaws}
  \lean{PhyXMiniProblems.ProblemPhyXMini0747.SatisfiesIdealProjectileLaws}
  \uses{def:physics:phyx-mini-0747:phyxminiproblems-problemphyxmini0747-timequantity, def:physics:phyx-mini-0747:phyxminiproblems-problemphyxmini0747-lengthinmeters, def:physics:phyx-mini-0747:phyxminiproblems-problemphyxmini0747-signedlengthinmeters, def:physics:phyx-mini-0747:phyxminiproblems-problemphyxmini0747-timeinseconds, def:physics:phyx-mini-0747:phyxminiproblems-problemphyxmini0747-signedvelocityinmeterspersecond, def:physics:phyx-mini-0747:phyxminiproblems-problemphyxmini0747-accelerationinmeterspersecondsquared, def:physics:phyx-mini-0747:phyxminiproblems-problemphyxmini0747-baseballwallsetup}
  Ideal no-drag projectile kinematics in coherent SI readouts. Time is elapsed from the hit, the horizontal origin is the hit point, and upward is positive. These general laws mention neither a wall-height answer nor an answer choice
\end{definition}
\begin{proof}
  This is the declared model definition of Satisfies Ideal Projectile Laws.
\end{proof}
\begin{definition}[Answer Choice]
  \label{def:physics:phyx-mini-0747:phyxminiproblems-problemphyxmini0747-answerchoice}
  \lean{PhyXMiniProblems.ProblemPhyXMini0747.AnswerChoice}
  Labels of the four wall-height answers printed with the problem
\end{definition}
\begin{proof}
  This is the declared model definition of Answer Choice.
\end{proof}
\begin{definition}[displayed Wall Height In Meters]
  \label{def:physics:phyx-mini-0747:phyxminiproblems-problemphyxmini0747-displayedwallheightinmeters}
  \lean{PhyXMiniProblems.ProblemPhyXMini0747.displayedWallHeightInMeters}
  \uses{def:physics:phyx-mini-0747:phyxminiproblems-problemphyxmini0747-answerchoice}
  Wall height in metres printed beside each displayed answer label
\end{definition}
\begin{proof}
  This is the declared model definition of displayed Wall Height In Meters.
\end{proof}
\begin{definition}[recorded Answer Choice]
  \label{def:physics:phyx-mini-0747:phyxminiproblems-problemphyxmini0747-recordedanswerchoice}
  \lean{PhyXMiniProblems.ProblemPhyXMini0747.recordedAnswerChoice}
  \uses{def:physics:phyx-mini-0747:phyxminiproblems-problemphyxmini0747-answerchoice}
  The answer label recorded in the supplied dataset metadata
\end{definition}
\begin{proof}
  This is the declared model definition of recorded Answer Choice.
\end{proof}
\begin{definition}[Matches Displayed Wall Height]
  \label{def:physics:phyx-mini-0747:phyxminiproblems-problemphyxmini0747-matchesdisplayedwallheight}
  \lean{PhyXMiniProblems.ProblemPhyXMini0747.MatchesDisplayedWallHeight}
  \uses{def:physics:phyx-mini-0747:phyxminiproblems-problemphyxmini0747-lengthinmeters, def:physics:phyx-mini-0747:phyxminiproblems-problemphyxmini0747-baseballwallsetup, def:physics:phyx-mini-0747:phyxminiproblems-problemphyxmini0747-answerchoice, def:physics:phyx-mini-0747:phyxminiproblems-problemphyxmini0747-displayedwallheightinmeters}
  A setup's wall-height readout agrees with a displayed answer
\end{definition}
\begin{proof}
  This is the declared model definition of Matches Displayed Wall Height.
\end{proof}
% --- Archon declaration topology end ---
% --- Archon physics formalization source end ---
